\documentclass[11pt]{article}

% Math
\usepackage{amsmath, amssymb, amsfonts, mathtools}

% Page layout
\usepackage[margin=1in]{geometry}

% Optional useful packages
\usepackage{physics}      % \bra, \ket, \dv, \pdv, etc.
\usepackage{bm}           % bold math symbols
\usepackage{bbm}
\usepackage{enumitem}     % better lists
\usepackage{graphicx}     % Required for inserting images
\usepackage{subcaption}   % Required for the subfigure environment
\usepackage{hyperref}     % links
\hypersetup{
    colorlinks=true,       % false: boxed links; true: colored links
    linkcolor=blue,        % color of internal links (sections, pages, etc.)
    citecolor=green,       % color of citation links to bibliography
    % filecolor=magenta,     % color of file links
    urlcolor=blue,         % color of external website links
    % allcolors=blue       % uncomment this if you want ALL links to use one color
}

\title{Weekly Update}
\author{Reza Asakereh}
\date{Last Update: \today}

\begin{document}

\maketitle

\section{Main Findings}
\subsection{DB-Sorter}
\textbf{quick background.} DB-sorter uses the sorting characteristic of DB flow to find ground state or excited states of a hamiltonian $H$ using a diagonal matrix $D$ that has an engineered order in its diagonal elements (usually increasing). 
\\

\textbf{Findings:}
\begin{itemize}
    \item This week was mainly spent debugging the warm starts and choosing a better warm start than simple VQE. Finally \href{https://github.com/rasakereh/double-bracket-algo/commit/6cc0a4a5a2879f2dcabb9d4beb3fffece7fb32bc}{VarQITE was implemented} and  \href{https://github.com/rasakereh/double-bracket-algo/commit/3016f4ae45ddad825df3e6b7bd4007539db1cbdc}{consistently used} for better results. The result was only a sharper drop at the first step while the rest of the steps wasn't any better.
    \item With some trials and errors, I can see that the base in which $D$ is diagonal has the boldest effect on the energy drop and convergence speed. This introduces a better place for "warm start"s. We can use a variational method as a heuristic to design a proper base for $D$. See \hyperref[sec:warm-start]{3.1.1}
\end{itemize}


\section{Plan For Next Week(s)}
\subsection{DB-sorter}
\begin{itemize}
    \item Finding fit basis for $D$.
    \item Wrapping up DB-sorter.
\end{itemize}


\section{Details}
\subsection{DB-sorter}
\subsubsection{D basis as warm start}
\label{sec:warm-start}
Previously, we saw that the energy dip can rely on the basis of $D$. Looking at the initial equation $\frac{dH}{dt} = [[D, H], H]$ one can assume deciding a basis that maximizes the double bracket whould be the way to go. For easier implementation, we try to only maximize the norm of $[D, H]$ on $\ket{k}$ for $k$'th excited state. Which is
\begin{align*}
    [D, H]\ket{v_k} &= (DH - HD)\ket{v_k} \\
    &= (UMU^*H - HUMU^*)U\ket{k} \\
    &= UMU^*HU\ket{k} - HUM\ket{k} \\
    &= UMU^*HU\ket{k} - m_0HU\ket{k} \\
    &= \sum{a_iUMU^*P_iU\ket{k} - a_im_0P_iU\ket{k}}
\end{align*}
We should be able to heuristicly do sth with it. While I am not sure how to proceed from here, my first idea is to run $l$ parallel variational circuits, one for each Pauli term in the Hamiltonian, with shared parameters defining $U$.

\end{document}
